%% kapa: exam=1 round=28 year=2017 period=2 form=A subject=appraisal_law subject_ko=감정평가관계법규 seq=25 number=25
%% source: Q-Net 감정평가사 기출문제 (cst00309) artlSeq=5199603 file=extracted/1008/5199603/2017년 제28회 감정평가사 1차시험 문제지/2교시 A형.pdf
%% conversion: kordoc 4.24.51 -> tools/gen1.py (auto); review=proofread-sonnet; flags=none
\begin{question}{25}
건축법령상 건축물의 대지와 도로에 관한 설명으로 옳지 않은 것은? (단, 건축법상 적용제외 규정 및 건축협정에 따른 특례는 고려하지 않음)
\begin{choices}
  \item 건축물의 주변에 허가권자가 인정한 유원지가 있는 경우에는 건축물의 대지가 자동차전용도로가 아닌 도로에 2 미터 이상 접할 것이 요구되지 아니 한다.
  \item 연면적의 합계가 3천 제곱미터인 작물 재배사의 대지는 너비 6 미터 이상의 도로에 4 미터 이상 접할 것이 요구되지 아니 한다.
  \item 주민이 오랫동안 통행로로 이용하고 있는 사실상의 통로로서 해당 지방자치단체의 조례로 정하는 것인 경우의「건축법」상 도로는 이해관계인의 동의를 받지 아니하고 건축위원회의 심의를 거쳐 그 도로를 폐지할 수 있다.
  \item 면적 5천 제곱미터 미만인 대지에 공장을 건축하는 건축주는 대지에 조경 등의 조치를 하지 아니할 수 있다.
  \item 도로면으로부터 높이 4.5 미터 이하에 있는 창문은 열고 닫을 때 건축선의 수직면을 넘지 아니하는 구조로 하여야 한다.
\end{choices}
\end{question}
\answer{3}
